% Folder 84, batch 08 (archivists' pages 141-160).
% Pass: Opus 5.5 (claude-opus-5-5), 2026-10-09 — first pass, unchecked against the pages by a human.
% Revised: Opus 5.5 (claude-opus-5-5), 2026-10-10 — re-read on the 700 dpi tiles; 28 \ill and 11 \uncertain resolved, 70 remain (34 \ill, 36 \uncertain). Page 143 read $f_i$, written over $u_i$, where the first pass had $u_i$.
%
% The even pages 142-160 are typescript pages of Récoltes et Semailles
% (pp. 474-489 of that text: notes 108-111, « Yang enterre yin », the
% « retrouvailles ») used as scrap paper; they carry no mathematics and only
% typing corrections, and are absent. The odd pages are the handwritten
% calculations, fast drafting, much of the connective prose illegible.
\documentclass[11pt,a4paper]{article}
\input{../preamble/grothendieck}

\folder{84}
\batch{8}
\pages{141}{160}
\foldertitle{[Formes quadratiques] : notes manuscrites (s.d.).}
\dating{[à partir de 1982-vers 1986]}
\watermark{Édition de démonstration}

\begin{document}

\page{141}
\note{la page reprend un calcul commencé avant la batterie : $F_{11}$, $E_i$, $L_i$, $M_i$, $\pi_i$, $\psi_i$ et la constante $\chi$ sont introduits plus haut dans le dossier}
Si on a pu définir $F_{11}$, \ill{} \ill{} des deux \ill{}
connus, dans \uncertain{l'hom} $F$ est un \uncertain{hom} d'\uncertain{ext.} \ill{} \ill{} (\uncertain{induisant}
id sur $L_1\otimes L_2$ et sur $M_1\otimes M_2$), donc est un iso.

\uncertain{Construction} \ill{} $F_{11}\colon E_1\otimes E_2\to E_1\otimes E_2$ \ill{}.

1°) Connu sur $E_1\otimes_0 E_2$, alors la \uncertain{composée}\note{mot en partie surchargé}
$E_1\otimes_0 E_2\xrightarrow{\ \pi\ } L_1\otimes L_2\hookrightarrow E_1\otimes E_2$
($\chi$-rétraction de $E_1\otimes_0 E_2$ sur le \uncertain{sous-module} $L_1\otimes L_2$),
i.e. $F_{11}\circ\mu$ est connu :

\begin{tikzcd}[column sep=large]
E_1\otimes_0 E_2 \arrow[r, hook, "\mu"] \arrow[dr, dashed, "\pi"'] & E_1\otimes E_2 \arrow[r, "F_{11}"] & E_1\otimes E_2 \\
 & L_1\otimes L_2 \arrow[ur, dashed, "\alpha_1\otimes\alpha_2"'] &
\end{tikzcd}

2°) $\nu\circ F_{11}$ est connu, de façon précise :

\begin{tikzcd}[column sep=large]
E_1\otimes E_2 \arrow[r, "F_{11}"] \arrow[dr, "\psi_1\otimes\psi_2"'] & E_1\otimes E_2 \arrow[r, "\nu"] & E_1\otimes^{0} E_2 \\
 & M_1\otimes M_2 \arrow[ur, "\varpi"'] &
\end{tikzcd}

\marginal{\struck{NB $\pi(\alpha_1\otimes L_2)=\chi\,\mathrm{id}_{L_1}\otimes L_2$}\quad
NB $\pi\circ(\alpha_1\otimes\alpha_2)=\chi\,\mathrm{id}_{L_1\otimes L_2}$,\quad
$(\psi_1\otimes\psi_2)\circ\varpi=\chi\,\mathrm{id}_{M_1\otimes M_2}$}

\[
\boxed{\begin{gathered}
\chi\,\mathrm{id}-(\pi'_1+\pi'_2)+\gamma\,\pi'_1\otimes\pi'_2\\
\gamma\chi=2 \qquad \text{e.g. } \chi=2,\ \gamma=1
\end{gathered}}
\]

dans le cas $\chi=2$, on \uncertain{aura}
\[
\left\{\begin{aligned}
F_{11}(\text{\struck{\ill{}}}\,x)&=\bar x+x\\
\bar x_i&=x_i-\pi_i(x_i)
\end{aligned}\right.
\qquad
\bar x_1\otimes\bar x_2=\bar x_1\otimes\bar x_2
\]
\note{la seconde formule est lue telle quelle ; le second membre est en partie effacé}
i.e. $\boxed{\pi_i(x_i)=x_i-\bar x_i}$

NB \uncertain{La} \uncertain{donnée} du $\pi_1$ équivaut à celle de $\sigma_1$
\struck{\ill{}} \uncertain{induisant} $-\mathrm{id}_{L_1}$ et $+\mathrm{id}_{M_1}$.

\page{143}
$E_1\circledast E_2$\note{le signe au-dessus de $E_2$ est surchargé et illisible}
\qquad
$E_1\circledast E_2=E_1\otimes E_2\oplus L_1\otimes L_2$, où $L_1\otimes L_2$ est accolé à
$(E_1\otimes_0 E_2)$

\note{dans les premières lignes de la page, $f$ surcharge partout un $u$ antérieur ($u_i$ devient $f_i$) ; seul ce qui suit « tq. » garde le $u$}
$f_i$ endom. de $E_i$ tq. $u_i(L_i)\subset L_i$, \quad $u_i|L_i=u_i^{!}\colon L_i\to L_i$.
\struck{$u_2$ endom. de $E_2$}
\begin{align*}
&f_1\otimes f_2 \text{ sur } E_1\otimes E_2\\
&f_1^{!}\otimes f_2^{!} \text{ sur } L_1\otimes L_2
\end{align*}
\struck{$u_1\otimes u_2$}
\begin{align*}
f(x_1\otimes x_2)&=f_1(x_1)\otimes f_2(x_2)\\
f(u_1\otimes u_2)&=f_1^{!}(u_1)\otimes f_2^{!}(u_2)
\end{align*}

\begin{tikzcd}[column sep=small, row sep=small]
 & E_1\otimes E_2 & \\
x_1\otimes u_2 \arrow[r] \arrow[dr] & x_1\otimes u_2 \arrow[r] & f_1(x_1)\circledast f_2(u_2) \\
 & \pi_1(x_1)\otimes u_2 \arrow[r] & f_1^{!}(\pi_1(x_1))\otimes f_2^{!}(u_2)
\end{tikzcd}

\note{les flèches du diagramme sont des $\mapsto$ sur la page ; à gauche, au-dessus de $x_1\otimes u_2$ : $E_1\otimes L_2$, avec le signe $\in$ couché}
\[
f_1(x_1)\circledast u_2=f_1^{!}(\pi_1(x_1))\otimes u_2
\]
\struck{$f(x$.}\quad \emph{Ex} $f_1=$ \quad \struck{\ill{}}
\[
f_1(x_1)=f_1(\pi\ \ldots
\]
\note{la formule s'interrompt sur le début d'un $\pi$}
\[
x_1-\pi_1(x_1)
\]
\[
E_1\otimes E_2\longrightarrow L_1\otimes L_2
\]
\note{sous la flèche, un $\alpha$ isolé ; le $L_1$ surcharge un $E_1$}

\begin{tikzcd}[column sep=small]
 & E_1\otimes E_2 \arrow[r] & L_1\otimes L_2 \\
E_1\otimes L_2 \arrow[ur] \arrow[r] & L_1\otimes L_2 \arrow[ur] &
\end{tikzcd}

\note{au-dessus de la flèche horizontale supérieure, une étiquette $\chi\, f_1\otimes f_2$ surchargée de traits verticaux (peut-être des « ! ») ; à droite de la ligne inférieure, une étiquette biffée, suivie de ce qui est peut-être $f_1^{!}\otimes f_2^{!}$}

\[
\begin{array}{cc}
E_1 & E_2\\
{\scriptstyle f_1}\downarrow & \downarrow{\scriptstyle f_2}\\
L_1 & L_2\\
\lambda_1 & \lambda_2
\end{array}
\qquad
\begin{array}{cc}
f_1 & f_2\\
\lambda_1 & \lambda_2
\end{array}
\qquad
\frac{\lambda_1}{\lambda_2}\,f_2
\]
\[
\chi\, f_1(x_1)\otimes \overbrace{f_2(u_2)}^{\lambda_2 u_2}
\]
\note{l'argument de $f_2$ surcharge un $x_2$ ; le $\lambda_2$ de l'accolade est repassé}
\struck{$f_1(f_1(x_1))\otimes u_2$}

$\lambda_1$ \struck{$\lambda_2$} $\pi_1(x_1)\otimes$ \ill{}

\struck{$f_1($} \quad $x_1\otimes\lambda_2 u_2$ \quad $\chi(\pi_1(x_1)\otimes u_2)$

\struck{$f_1\otimes f_2$} \quad \struck{$\lambda_1\lambda_2(\pi_1(x_1)\otimes u_2)$}
\quad $f_1(\pi_1(x_1))\otimes f_2(u_2)$

\struck{$\lambda_1 f_1(x_1)=\lambda$}
\[
\lambda_1\pi_1(x_1)=f_1(x_1)
\]

\page{145}
\[
T_1=2\cdot 1_1,\,-b_1\ =\ \begin{pmatrix}2&0\\-b_1&0\end{pmatrix}
\qquad
\bar T_1=\begin{pmatrix}0&0\\b_1&2\end{pmatrix}
\]
\note{le « $2$ » de $T_1$ surcharge un autre signe ; dans $\bar T_1$ les deux zéros de la première ligne sont cerclés}

$\boxed{b=b_1b_2}$\note{cerclé} \qquad \struck{\ill{}}

$(E_1,\xi_1),(E_2,\xi_2)\longrightarrow(E,\xi)$ \uncertain{avec}\note{mot repassé, peut-être biffé}, $\downarrow$ triv.
\qquad\struck{\ill{}}\qquad $b_1$ \quad $b_2$

\struck{$b_1=E_1$}
\begin{align*}
b_1 &: M_1\xrightarrow{\ \tau_1\ }E_1\xrightarrow{\ \pi_1\ }L_1\\
-b_1 &: M_1\xrightarrow{\ \varpi_1\ }E_1\xrightarrow{\ \tau_1\cdot\ }L_1
\end{align*}
trivialisations $\xi\simeq\beta$
(\ill{}-trivialisations)
\[
\left.\begin{aligned}
\tau'_1&=\tau_1+u_1\\
b'_1&=b_1+\chi u_1
\end{aligned}\ \right|\
\begin{aligned}
\tau'_2&=\tau_2+u_2\\
b'_2&=b_2+\chi u_2
\end{aligned}
\qquad u_i\in\mathrm{Hom}(M_i,L_i)
\]
\[
\beta_*(\xi)=b\in\mathrm{Hom}(M_1,L_1)\qquad \beta'(\eta)=-b
\]
\note{devant $\beta_*$, un mot biffé}
\begin{align*}
x_1&\mapsto \pi_1(\tau_1^{*}(x_1))\\
x_1&\mapsto \pi_1(\tau_1(x_1)+u_1(x_1))=\pi_1(\tau_1(x_1))
\end{align*}

\page{147}
\begin{align*}
\sigma_1&=\mathrm{id}_1-\gamma\,\pi_1\\
\sigma_2&=\mathrm{id}_2-\gamma\,\pi_2
\end{align*}
\begin{align*}
\mathrm{id}_{12}+\sigma_1\otimes\sigma_2
&=\mathrm{id}_{12}+\mathrm{id}_{12}-\gamma(\pi'_1+\pi'_2)+\gamma^2\pi'_1\otimes\pi'_2\\
&=2\,\mathrm{id}_{12}-\gamma(\pi'_1+\pi'_2)+\gamma^2\pi'_1\otimes\pi'_2
\end{align*}
\note{le second $\mathrm{id}_{12}+{}$ est ajouté au-dessus de la ligne et relié par un trait}
\[
\boxed{\gamma\chi=2}
\qquad
\begin{aligned}
\beta_1&=\beta_2=1-\gamma\chi\\
\alpha&=-\chi(1-\gamma\chi)
\end{aligned}
\qquad
\boxed{\begin{aligned}\beta&=-1\\ \alpha&=\chi\end{aligned}}
\]

\struck{$\sigma_1=\alpha_1+\beta_1\pi_1$}
\[
\sigma_i=\lambda 1_i+\mu\pi_i
\]
\begin{align*}
\mathrm{id}+\sigma_1\otimes\sigma_2
&=\underbrace{\mathrm{id}_{12}+\lambda^2\mathrm{id}_{12}}_{(1+\lambda^2)\,\mathrm{id}_{12}}
+\lambda\mu(\pi'_1+\pi'_2)+\mu^2\pi'_1\otimes\pi'_2
\end{align*}
\[
\boxed{\left\{\begin{aligned}
\lambda^2+1&=\chi\\
\lambda\mu&=-1\\
\mu^2&=\gamma
\end{aligned}\right.}
\]

$u_i$ \uncertain{\ill{}} de \struck{\ill{}} $E_i$ \uncertain{i.e.} \ill{} \ill{}

$\mathcal{H}_i=$ \ill{} \ill{} \qquad $1,0$

\[
0\to L\to E\to M\to 0
\]
\[
0\to\mathrm{Hom}(M,L)\longrightarrow\mathcal{H}/\mathrm{id}\longrightarrow\frac{\mathcal{O}\times\mathcal{O}}{\mathcal{O}}\to 0
\qquad
\to\mathcal{H}\to\mathcal{O}\to
\]
\note{la flèche de gauche porte une étiquette biffée ; $\mathcal{H}$ est barré d'un trait oblique suivi de « id » en indice, lu ici comme un quotient}
\begin{align*}
\sigma(u)&=(\lambda(u)+\mu(u))\,\mathrm{id}_E-u\\
u+\sigma(u)&=(\lambda(u)+\mu(u))\,\mathrm{id}
\end{align*}
\[
\begin{array}{c}
\mathrm{Hom}(M\otimes M',\,L\otimes L')\\
\uparrow\\
\mathrm{Hom}(M,L)\otimes\mathrm{Hom}(M',L')
\end{array}
\]

\page{149}
\struck{Mais \uncertain{je} \uncertain{veux} compatibilité \uncertain{avec} les $\pi$-trivialisations.}
\note{la première ligne et le petit diagramme qui suit sont barrés de traits obliques}

\struck{$E_1\vee E_2\ \big\downarrow\,\pi_{E_1\vee E_2}$}
\[
\bigl(x_1\otimes x_2+\bar x_1\otimes\bar x_2\bigr)
\quad
\bigl(x_1-\pi_1(x_1)\bigr)\bigl(x_2-\pi_2(x_2)\bigr)\bigr)
\]
\note{la première expression est entourée ; la seconde, sous $\bar x_1\otimes\bar x_2$, en développe les facteurs}
\begin{align*}
\bar{\bar x}_1&=\bigl(x_1-\pi_1(x_1)\bigr)-\underbrace{\pi_1\bigl(x_1-\pi_1(x_1)\bigr)}_{-\pi_1(x_1)+\chi\pi_1(x_1)}=x_1\\
&\hphantom{=}\qquad\text{avec } \chi=2
\end{align*}
\note{au-dessus de la parenthèse, « $x_1$, $0$ » biffé}

$x_1\mapsto\bar x_1$ induit $-\mathrm{id}$ sur $L_1$, $+\mathrm{id}$ sur $M_1=E_1/L_1$
\[
\pi^2-\chi\pi=0\qquad \pi(\pi-\chi)=0 \qquad 0,\ \chi
\]
\[
\varpi_1(x_1)=2x_1-\pi_1(x_1)
\]
\begin{align*}
x_1-\varpi_1(x_1)
&=x_1-\bigl(2x_1-\pi_1(x_1)\bigr)=\pi_1(x_1)-x_1=-\bar x_1
\end{align*}
\[
\underbrace{(\chi-1)\,\mathrm{id}+(1-\pi_1)\otimes(1-\pi_2)}+(\gamma-1)\,\pi_1\otimes\pi_2
\]
\note{sous les deux « $1$ » de $(1-\pi_1)\otimes(1-\pi_2)$, un petit $1$ en indice}

\page{151}
\[
\underbrace{\beta_1\doteq\beta_2}_{\beta}=1-\gamma\chi
\qquad
\alpha=-\beta\chi=-\chi(1-\gamma\chi)
\]
\note{sous les deux occurrences de $\gamma$, trois petits traits horizontaux}
\[
-\chi(1-\gamma\chi)\,\mathrm{id}+(1-\gamma\chi)(\pi'_1+\pi'_2)+\gamma\,\pi'_1\otimes\pi'_2
\qquad \gamma\in\mathbb{Z}
\]
\[
-\chi\underbrace{(1-\gamma\chi)}+2(1-\gamma\chi)\chi+\gamma\chi^2=\chi
\qquad
\text{\struck{$2\chi-\chi$}}\ \chi=\chi \quad \text{o.k.!}
\]
\struck{$\boxed{2\neq\chi}$}
\begin{align*}
&-\chi\cdot(1-\gamma\chi)\,\mathrm{id}_{M_1\otimes M_2}
+(1-\gamma\chi)\bigl[(\tau_1\otimes\mathrm{id}_2)+\mathrm{id}_1\otimes\tau_2\bigr]\\
&\quad\text{\struck{$x_1\otimes x_2$}}+\gamma\,\tau_1\otimes\tau_2
\end{align*}
\[
\left\{\begin{aligned}
-\chi(1-\gamma\chi)&=\chi\\
1-\gamma\chi&=-1
\end{aligned}\right.
\qquad
1-\gamma\chi=-1
\qquad
\boxed{\gamma\chi=2}
\]
ok \uncertain{si} $\chi=\pm 2$ \struck{\ill{}}, \uncertain{si} $\chi$\note{la ligne s'arrête sur $\chi$}

$\chi=2$, $\gamma=1$, $\beta=-1$, $\alpha=2$, donc \uncertain{on} prend
\[
\text{\struck{$2\,\mathrm{id}$}}\quad
2\,\mathrm{id}_{E_1\otimes E_2}-(\pi_{E_1}\otimes\mathrm{id}_{E_2}+\mathrm{id}_{E_1}\otimes\pi_{E_2})
+\pi_{E_1}\otimes\pi_{E_2}
\]
\note{une flèche courbe renvoie de « $\beta=-1$, $\alpha=2$ » à cette formule}
\begin{align*}
x_1\otimes x_2\longmapsto{}& 2x_1\otimes x_2-\pi_1(x_1)\otimes x_2-x_1\otimes\pi_2(x_2)\\
&+\pi_1(x_1)\otimes\pi_2(x_2)
\end{align*}
\struck{$\equiv x_1\otimes x_2+(\bar x_1\otimes\bar x_2)$}
\begin{align*}
\bar x_1&=\tau_1\psi_1(x_1)-x_1=x_1-\pi_1(x_1)\\
\pi_1(x_1)&=2x_1-\tau_1\psi_1(x_1)
\end{align*}

\page{153}
\[
\text{\struck{$E_1\otimes E_2\to$}}\quad
L_1\otimes L_2+M_1\otimes L_2+L_1\otimes M_2+M_1\otimes M_2
\]
\note{deux accolades sous la ligne regroupent les termes deux à deux ; dans le deuxième terme, $M_1$ surcharge un $L_1$ et le second facteur, repassé, est lu $L_2$}

\begin{tikzcd}[column sep=large]
E_1\otimes_0 E_2 \arrow[r, "F_{11}^{!}"] \arrow[d] \arrow[dr, "\pi"] & E_1\otimes E_2 \arrow[r] & E_1\otimes^0 E_2 \\
E_1\otimes E_2 \arrow[dr, "\psi_1\otimes\psi_2"'] & L_1\otimes L_2 \arrow[u, "\alpha_1\otimes\alpha_2"'] & \\
 & M_1\otimes M_2 \arrow[uur, bend right=20] &
\end{tikzcd}

\note{à côté de $\alpha_1\otimes\alpha_2$, un mot souligné, peut-être « \uncertain{connu} » ; une flèche courbe relie ce mot à la flèche de $M_1\otimes M_2$ vers $E_1\otimes^0 E_2$ ; le placement exact des flèches inférieures est simplifié}
\qquad $(E_1\otimes M_2)\times(M_1\otimes E_2)$

$L_1\otimes L_2$ --- $\otimes$ --- $\otimes$\note{deux $\otimes$ cerclés}

\[
\begin{array}{c|cccc}
 & L_1\otimes L_2 & L_1\otimes M_2 & M_1\otimes L_2 & M_1\otimes M_2\\
\hline
L_1\otimes L_2 & \chi & \mathrm{id}_{L_1}\otimes\tau_2 & \tau_1\otimes\mathrm{id}_{L_2} & ?\ (\tau_1\otimes\tau_2\,?)\\
L_1\otimes M_2 & 0 & 0 & 0 & \doteq\ \tau_1\otimes\mathrm{id}_{M_2}\\
M_1\otimes L_2 & 0 & 0 & 0 & \doteq\ \mathrm{id}_{M_1}\otimes\tau_2\\
M_1\otimes M_1 & 0 & 0 & 0 & \chi
\end{array}
\]
\note{« $M_1\otimes M_1$ » sic en tête de la dernière ligne ; une flèche verticale le long de la colonne d'en-tête remonte de la dernière ligne vers la première}

On trouve que les $2$ \struck{\ill{} \ill{}} \ill{} \uncertain{connues}\add{sur $F_{11}$}
sont compatibles, et laissent une indéterminée, qui est un hom de $M_1\otimes M_2$ dans
$L_1\otimes L_2$ :
\[
E_1\otimes E_2\qquad E_1\otimes E_2
\]
\[
\pi'_1=\alpha_1\pi_1\ \Big|\ \varpi'_1=\varpi_1\psi_1\qquad \pi'_1+\varpi
\]
\[
\pi'_2=\alpha_2\pi_2\qquad \varpi'_2=\omega_2\psi_2
\]

C'est normal, on \uncertain{pourrait} l'\ill{} \uncertain{par} \ill{} \uncertain{que}
le fait que \ill{} \ill{} \ill{} \uncertain{d'extensions} --- et celle-ci est \ill{},
on a l'\uncertain{indétermination} des choix !

\page{155}
\struck{$-\chi(1-\delta\chi)$}
\[
\boxed{\begin{aligned}
&-\chi(1-\delta\chi)\,x_1\otimes x_2+(1-\delta\chi)\bigl[x_1\otimes\pi_2(x_2)+\pi_1(x_1)\otimes x_2\bigr]\\
&\qquad+\delta\,\pi_1(x_1)\otimes\pi_2(x_2)
\end{aligned}}
\]
\note{au-dessus de l'encadré, une lettre isolée, peut-être $\xi$ ; $\delta$ joue ici le rôle de la constante notée $\gamma$ aux pages précédentes}

Si $\chi=2$, on peut prendre $\delta=1$, donc
\begin{align*}
\alpha&=-2(1-2)=2\\
\beta&=\gamma=1-2=-1
\end{align*}
donc
\[
2x_1\otimes x_2-\bigl[(x_1\otimes\pi_2(x_2)+\pi_1(x_1)\otimes x_2)\bigr]+\pi_1(x_1)\otimes\pi_2(x_2)
\]
\begin{align*}
&\text{\struck{$\pi_i(x_i)=2x_i-\tau_i\psi_i(x_i)$}}\\
\pi_i(x_i)&=2x_i-\tau_i\psi_i(x_i)=2x_i-(x_i+\bar x_i)=x_i-\bar x_i\\
\bar x_i&=\tau_i\psi_i(x_i)-x_i\\
x_i+\bar x_i&=\tau_i\psi_i(x_i)
\end{align*}
\begin{align*}
&2x_1\otimes x_2-\bigl[(x_1\otimes(x_2-\bar x_2))+(x_1-\bar x_1)\otimes x_2\bigr]\\
&\qquad+(x_1-\bar x_1)\otimes(x_2-\bar x_2)\\
&=x_1\otimes x_2+\bar x_1\otimes\bar x_2 \qquad\text{o.k.}
\end{align*}
\begin{align*}
\bar x_i&=x_i-\pi_i(x_i)=\Bigl\{\\
&x_1\otimes x_2+\text{\struck{$(x_i-\pi_i(x_i))$}}
\end{align*}
\[
-\chi(1-\delta\chi)\overset{?}{=}\chi\qquad -1+\delta\chi=1\qquad \boxed{\delta\chi=2}
\]

\page{157}
\struck{$\alpha\,u_1\otimes x_2+$}\quad\struck{$\alpha\,\ill{}+\chi\gamma$}
\[
(\alpha+\chi\gamma)\,u_1\otimes x_2+(\beta+\delta\chi)\,u_1\otimes\pi_2(x_2)
\]
\note{sous le second terme, « $\shortparallel$ $u_1\otimes\pi_2(x_2)$ » ; $\beta$ y surcharge une lettre}
\[
(\alpha+\chi\gamma)\,u_1\otimes x_2+(\beta+\delta\chi-1)\,u_1\otimes\pi_2(x_2)=0
\]
$\forall u_1\in L_1$, $x_2\in E_2$,

si $x_2=u_2$, \uncertain{on} \uncertain{trouve}
\[
(\alpha+\chi\gamma)+\chi(\beta+\delta\chi-1)=0
\]
\[
u_1\otimes\bigl[(\alpha+\chi\gamma)\,x_2+(\beta+\delta\chi-1)\,\pi_2(x_2)\bigr]=0
\]
\note{au-dessus de la première de ces deux formules, un « $u_1$ » isolé}

Donc dans $L^{!}_2$ \uncertain{on} \uncertain{a} :
\[
(\alpha+\chi\gamma)\,\mathrm{id}=0
\qquad
\alpha+\chi\gamma=0
\qquad
\alpha=-\chi\gamma
\]
donc \ill{} $\alpha=-\chi\beta$ \quad donc $\underline{\beta=\gamma}$

\struck{$\beta+\delta$}
\[
-\chi\beta\,x_1\otimes x_2+\beta\bigl(x_1\otimes\pi_2(x_2)+\pi_1(x_1)\otimes x_2\bigr)+\delta\,\pi_1(x_1)\otimes\pi_2(x_2)
\]
il \uncertain{reste} \uncertain{donc} \struck{\ill{}}
\[
\text{\struck{$\ldots\oplus(\beta\ldots$}}\quad(\beta+\delta\chi-1)\,\pi_2(x_2)=0
\]
\[
\boxed{\beta+\delta\chi=1}
\qquad
\begin{aligned}
\beta&=1-\delta\chi\\
\alpha&=-\chi\beta=\text{\struck{$-\chi+\chi^2$}}
\end{aligned}
\]
\[
\beta=\beta=1-\delta\chi\qquad \alpha=-\chi(1-\delta\chi)
\]
\note{la dernière ligne est écrite en bas à gauche, sous le numéro du feuillet ; le second $\beta$ surcharge une autre lettre}

\page{159}
\[
\begin{gathered}
0\to L_1\to E_1\to L'_1\to 0\\
0\to L_2\to E_2\to L'_2\to 0
\end{gathered}
\]
$L_1$, $L'_1$ \uncertain{facteurs} \ill{}. Mais \uncertain{supposons} $L'_1$, $L'_2$ loc.\ libres de rang fini.

\note{dans les deux suites exactes, $L_1$ et $L_2$ surchargent $E$ ; dans la seconde, et devant $\otimes L_2$ à la ligne suivante, $L$ porte un exposant qui a la forme d'un « 4 », transcrit « ! » par analogie avec les pages précédentes}
\[
\longrightarrow\qquad 0\to L^{!}_1\otimes L_2\to E_1\wedge E_2\to L'_1\otimes L'_2\to 0
\]
\[
0\leftarrow L_1^{\vee}\otimes L_2^{\vee}\leftarrow E_1^{\vee}\wedge E_2^{\vee}\leftarrow L_1'^{\vee}\otimes L_2'^{\vee}\leftarrow 0
\]
\note{des accolades sous la seconde suite relient $L_1^{\vee}\otimes L_2^{\vee}$ et $L_1'^{\vee}\otimes L_2'^{\vee}$ aux termes de la suite précédente}
$\exists$ des \uncertain{qu.} \quad $E_1^{\vee}\wedge E_2^{\vee}\simeq(E_1\wedge E_2)^{\vee}$
\[
\mathrm{Hom}(L'_1,\mathcal{O})\otimes\mathrm{Hom}(L'_2,\mathcal{O})\simeq\mathrm{Hom}(L'_1\otimes L'_2,\mathcal{O})
\]
\[
\text{\struck{$f_1$}}\otimes f_2+\bar f_1\otimes\bar f_2,\qquad u_1\otimes u_2
\]
\[
\begin{array}{c}
A_1\otimes A_2 \curvearrowleft\\
\downarrow\quad A_1\otimes_0 A_2\\
L_1\otimes L_2\quad \pi_1\otimes\mathrm{id}\\
\phantom{L_1\otimes L_2}\quad\mathrm{id}\otimes\pi_2
\end{array}
\qquad
\begin{aligned}
&\bigl(\mathrm{Tr}_1(f_1)-f_1\bigr)\otimes\bigl(\mathrm{Tr}_2(f_2)-f_2\bigr)\\
&\bigl(\tau'_1(f_1)-f_1\bigr)\bigl(\pi_2\psi_2)-f_2\bigr)
\end{aligned}
\]
\note{à gauche, une flèche verticale de $A_1\otimes A_2$ vers $L_1\otimes L_2$ et une flèche oblique remontant vers $A_1\otimes_0 A_2$ ; la seconde ligne de droite est lue telle quelle, les parenthèses n'y sont pas équilibrées}
\[
\text{\struck{$f_1$}}\qquad \text{\struck{$f_2$}}\qquad f_1
\]
\[
\pi_1(x_1)-\ldots
\qquad
\text{\struck{$\sigma=\pi$}}\qquad
\sigma(x)=\pi(x)\cdot 1-x
\qquad
\sigma=\pi\otimes e-\mathrm{id}
\]
\note{dans $\pi\otimes e$, la lettre lue $e$ est peut-être un $\ell$}
\[
\begin{aligned}
\pi_1(x_1)\qquad x_1\otimes x_2&\longmapsto \text{\struck{$\pi_1(x_1)\otimes x$}}\quad u_1\otimes\pi_2(x_2)\\
x_1\otimes x_2&\longmapsto \pi_1(x_1)\otimes u_2
\end{aligned}
\qquad
u_1\otimes\underline{\pi\sigma(x_2)}=\ldots
\]
\note{le second membre de la dernière formule s'interrompt après le signe d'égalité}
\[
\boxed{x_1\otimes x_2+\pi'_1(x_1)\otimes\pi'_2(x_2)}
\]
\[
\left\{\begin{aligned}
&\pi_2(x)=\overbrace{x}^{(\chi-1)x}+\pi'_2(x)\\
&\text{si } x\in L_2,\ \ldots\\
&\pi'_2(x)=x
\end{aligned}\right.
\]
\marginal{écrit en oblique : $2x-\psi(x)\tau$ ; $x-\psi(x)\tau=-\sigma(x)$}
\note{la page s'achève sur ce calcul ; rien ne montre où il se poursuit}

\end{document}
